% \iffalse meta-comment
%
% hit-paragraph.dtx 是段落与行距
%
% \fi
%
% \section{段落与行距}
%
% ^^A ======================================================================
% ^^A Module thesis-paragraph: 段落首行缩进与标点处理（hit-thesis）
% ^^A ======================================================================
%    \begin{macrocode}
%<*thesis-paragraph>
\ctexset{%
  punct=quanjiao,
  space=auto,
  autoindent=true}
%    \end{macrocode}
%
% 调整默认列表环境间的距离，以符合中文习惯。机制在
% \file{src/utils/hit-list.dtx}，这里只给取值。
%    \begin{macrocode}
\bool_gset_true:N \g__hit_list_nosep_bool
\clist_gset:Nn \g__hit_list_itemize_clist
  { leftmargin=0em,itemsep=0em,partopsep=0em,parsep=0em,topsep=0em,itemindent=3em }
\clist_gset:Nn \g__hit_list_enumerate_clist
  { leftmargin=0em,itemsep=0em,partopsep=0em,parsep=0em,topsep=0em,itemindent=3.5em }
\hit@setup@list@spacing
%    \end{macrocode}
%
% \changes{v3.2a}{0000/00/00}{列表间距的三行迁进 \file{src/utils/hit-list.dtx}，
%   与 \env{publist} 同处；这里只留调用与正文首行缩进}
% 正文首行缩进。
%    \begin{macrocode}
\dim_set:Nn \parindent {2em}
%    \end{macrocode}
%
% \changes{v3.2a}{0000/00/00}{\cs{cmd}/\cs{cs} 排版命令名工具从 book-math 迁入 book-paragraph（归位至正文排版细节）}
% 排版命令名的辅助：\cs{cs} 在文本中打印 \cs{name} 字面。
%    \begin{macrocode}
\cs_set:Npn \cmd #1 { \cs { \exp_after:wN \cmd@to@cs \token_to_str:N #1 } }
\cs_set:Npn \cmd@to@cs #1#2 { \char \number `#2 \relax }
\DeclareRobustCommand\cs[1]{\texttt{\char`\\#1}}
%</thesis-paragraph>
%    \end{macrocode}
%
% \Finale
